\documentclass[a4paper,12pt]{book}

% === Encodage et langue ===
\usepackage[utf8]{inputenc}
\usepackage[T1]{fontenc}
\usepackage[french]{babel}

% === Mathématiques ===
\usepackage{amsmath}
\usepackage{amssymb}
\usepackage{amsthm}
\usepackage{mathtools}
\usepackage{stmaryrd}

% === Boîtes colorées ===
\usepackage[theorems,skins,breakable]{tcolorbox}

% === Diagrammes et graphiques ===
\usepackage{tikz}
\usepackage{tikz-cd}
\usepackage{pgfplots}
\pgfplotsset{compat=1.18}

% === Figures et tableaux ===
\usepackage{graphicx}
\usepackage{float}
\usepackage{caption}
\usepackage{booktabs}
\usepackage{array}
\usepackage{longtable}

% === Listes ===
\usepackage{enumitem}

% === Références et index ===
\usepackage{hyperref}
\usepackage{cleveref}
\usepackage{makeidx}
\usepackage{minitoc}

\makeindex

% ============================================================
% Opérateurs mathématiques
% ============================================================
\DeclareMathOperator{\Hom}{Hom}
\DeclareMathOperator{\End}{End}
\DeclareMathOperator{\Aut}{Aut}
\DeclareMathOperator{\im}{im}
\DeclareMathOperator{\Id}{Id}
\DeclareMathOperator{\Ker}{Ker}
\DeclareMathOperator{\coker}{coker}
\DeclareMathOperator{\Ob}{Ob}
\DeclareMathOperator{\Mor}{Mor}
\DeclareMathOperator{\Fun}{Fun}
\DeclareMathOperator{\Nat}{Nat}
\DeclareMathOperator{\colim}{colim}
\DeclareMathOperator{\eq}{eq}
\DeclareMathOperator{\coeq}{coeq}

% Catégories usuelles
\DeclareMathOperator{\Set}{\mathbf{Set}}
\DeclareMathOperator{\Grp}{\mathbf{Grp}}
\DeclareMathOperator{\Ab}{\mathbf{Ab}}
\DeclareMathOperator{\Ring}{\mathbf{Ring}}
\DeclareMathOperator{\Top}{\mathbf{Top}}
\DeclareMathOperator{\Vect}{\mathbf{Vect}}
\DeclareMathOperator{\Mod}{\mathbf{Mod}}
\DeclareMathOperator{\Cat}{\mathbf{Cat}}

% Commandes utiles
\newcommand{\op}{\mathrm{op}}
\newcommand{\id}{\mathrm{id}}
\newcommand{\N}{\mathbb{N}}
\newcommand{\Z}{\mathbb{Z}}
\newcommand{\Q}{\mathbb{Q}}
\newcommand{\R}{\mathbb{R}}
\newcommand{\C}{\mathbb{C}}
\newcommand{\K}{\mathbb{K}}

% ============================================================
% Environnements de théorèmes
% ============================================================
\newtheorem{theoreme}{Théorème}[chapter]
\newtheorem{proposition}[theoreme]{Proposition}
\newtheorem{lemme}[theoreme]{Lemme}
\newtheorem{corollaire}[theoreme]{Corollaire}
\newtheorem{definition}[theoreme]{Définition}
\newtheorem{exemple}[theoreme]{Exemple}
\newtheorem{exercice}{Exercice}[chapter]
\newtheorem*{remarque}{Remarque}
\newtheorem*{notation}{Notation}

% === Styles tcolorbox ===
\tcolorboxenvironment{definition}{
  enhanced, breakable,
  colback=blue!5, colframe=blue!70!black,
  boxrule=0.8pt, arc=2pt,
  left=6pt, right=6pt, top=4pt, bottom=4pt
}
\tcolorboxenvironment{theoreme}{
  enhanced, breakable,
  colback=red!5, colframe=red!70!black,
  boxrule=0.8pt, arc=2pt,
  left=6pt, right=6pt, top=4pt, bottom=4pt
}
\tcolorboxenvironment{proposition}{
  enhanced, breakable,
  colback=red!3, colframe=red!60!black,
  boxrule=0.8pt, arc=2pt,
  left=6pt, right=6pt, top=4pt, bottom=4pt
}
\tcolorboxenvironment{lemme}{
  enhanced, breakable,
  colback=yellow!10, colframe=yellow!70!black,
  boxrule=0.8pt, arc=2pt,
  left=6pt, right=6pt, top=4pt, bottom=4pt
}
\tcolorboxenvironment{corollaire}{
  enhanced, breakable,
  colback=orange!5, colframe=orange!70!black,
  boxrule=0.8pt, arc=2pt,
  left=6pt, right=6pt, top=4pt, bottom=4pt
}
\tcolorboxenvironment{exemple}{
  enhanced, breakable,
  colback=green!5, colframe=green!60!black,
  boxrule=0.8pt, arc=2pt,
  left=6pt, right=6pt, top=4pt, bottom=4pt
}
\tcolorboxenvironment{exercice}{
  enhanced, breakable,
  colback=purple!5, colframe=purple!70!black,
  boxrule=0.8pt, arc=2pt,
  left=6pt, right=6pt, top=4pt, bottom=4pt
}
\tcolorboxenvironment{remarque}{
  enhanced, breakable,
  colback=gray!5, colframe=gray!60!black,
  boxrule=0.8pt, arc=2pt,
  left=6pt, right=6pt, top=4pt, bottom=4pt
}

% ============================================================
\title{\Huge\textbf{Théorie des Catégories} \\ \Large Cours de Master}
\author{Notes de cours}
\date{2025--2026}

\begin{document}

\maketitle
\dominitoc
\tableofcontents

% ============================================================
% PRÉFACE
% ============================================================
\chapter*{Préface}
\addcontentsline{toc}{chapter}{Préface}

La théorie des catégories\index{catégorie}, née des travaux fondateurs d'Eilenberg et Mac\,Lane dans les années 1940, s'est imposée comme le langage unificateur des mathématiques modernes. Conçue initialement pour formaliser la notion de \emph{transformation naturelle}\index{transformation naturelle} en topologie algébrique, elle a trouvé des applications profondes dans pratiquement toutes les branches des mathématiques.

C'est avant tout dans l'\oe uvre monumentale d'Alexandre Grothendieck\index{Grothendieck} que la théorie des catégories a révélé sa pleine puissance. En refondant la géométrie algébrique sur des bases catégoriques --- schémas, faisceaux, topos --- Grothendieck a démontré qu'un cadre abstrait bien choisi permet non seulement d'unifier des théories existantes, mais de découvrir des structures invisibles dans les formulations classiques.

Ce cours s'adresse aux étudiants de Master souhaitant acquérir une maîtrise solide des concepts fondamentaux : catégories, foncteurs, transformations naturelles, limites, adjonctions, et au-delà. L'accent est mis sur la rigueur des définitions et la richesse des exemples, issus de l'algèbre, de la topologie et de l'analyse fonctionnelle.

\bigskip

\chapter*{Notations}
\addcontentsline{toc}{chapter}{Notations}

\begin{notation}
Tout au long de ce texte, on adopte les conventions suivantes :
\begin{itemize}[leftmargin=2cm]
  \item[$\Set$] --- catégorie des ensembles et applications.
  \item[$\Grp$] --- catégorie des groupes et homomorphismes.
  \item[$\Ab$] --- catégorie des groupes abéliens et homomorphismes.
  \item[$\Ring$] --- catégorie des anneaux (unitaires) et homomorphismes d'anneaux.
  \item[$\Top$] --- catégorie des espaces topologiques et applications continues.
  \item[$\Vect_\K$] --- catégorie des $\K$-espaces vectoriels et applications linéaires.
  \item[$R\text{-}\Mod$] --- catégorie des $R$-modules à gauche et homomorphismes de modules.
  \item[$\Cat$] --- catégorie des petites catégories et foncteurs.
  \item[$\mathcal{C}^{\op}$] --- catégorie opposée\index{catégorie!opposée} de $\mathcal{C}$.
  \item[$\Hom_{\mathcal{C}}(A,B)$] --- ensemble des morphismes de $A$ vers $B$ dans $\mathcal{C}$.
  \item[$\id_A$] --- morphisme identité de l'objet $A$.
  \item[$F \circ G$] --- composition de foncteurs.
  \item[$\alpha : F \Rightarrow G$] --- transformation naturelle de $F$ vers $G$.
  \item[$[\mathcal{C}, \mathcal{D}]$] --- catégorie des foncteurs de $\mathcal{C}$ vers $\mathcal{D}$.
\end{itemize}
\end{notation}

% ============================================================
% CHAPITRE 1
% ============================================================
\chapter{Catégories, Foncteurs, Transformations Naturelles}
\minitoc

\section{Catégories}

\subsection{Définition d'une catégorie}

\begin{definition}[Catégorie]\label{def:categorie}
\index{catégorie!définition}
Une \textbf{catégorie} $\mathcal{C}$ est la donnée de :
\begin{enumerate}[label=(\roman*)]
  \item une classe $\Ob(\mathcal{C})$ dont les éléments sont appelés \textbf{objets}\index{objet} de $\mathcal{C}$ ;
  \item pour tout couple $(A, B)$ d'objets, un ensemble $\Hom_{\mathcal{C}}(A, B)$ dont les éléments sont appelés \textbf{morphismes}\index{morphisme} (ou \textbf{flèches}) de $A$ vers $B$ ;
  \item pour tout triplet $(A, B, C)$ d'objets, une loi de \textbf{composition}\index{composition}
  \[
    \circ : \Hom_{\mathcal{C}}(B, C) \times \Hom_{\mathcal{C}}(A, B) \longrightarrow \Hom_{\mathcal{C}}(A, C), \quad (g, f) \longmapsto g \circ f ;
  \]
  \item pour tout objet $A$, un morphisme $\id_A \in \Hom_{\mathcal{C}}(A, A)$ appelé \textbf{identité}\index{morphisme!identité} de $A$ ;
\end{enumerate}
satisfaisant les axiomes suivants :
\begin{enumerate}[label=(C\arabic*)]
  \item \textbf{Associativité}\index{associativité} : pour tous morphismes $f : A \to B$, $g : B \to C$, $h : C \to D$,
  \[ h \circ (g \circ f) = (h \circ g) \circ f. \]
  \item \textbf{Identité} : pour tout morphisme $f : A \to B$,
  \[ f \circ \id_A = f = \id_B \circ f. \]
\end{enumerate}
On exige de plus que les ensembles $\Hom_{\mathcal{C}}(A, B)$ soient deux à deux disjoints pour des couples $(A,B)$ distincts (condition de « typage » des morphismes).
\end{definition}

\begin{remarque}
Si $f \in \Hom_{\mathcal{C}}(A,B)$, on note aussi $f : A \to B$ ou $A \xrightarrow{f} B$. L'objet $A$ est la \textbf{source}\index{source} (ou domaine) de $f$ et $B$ est le \textbf{but}\index{but} (ou codomaine).
\end{remarque}

\subsection{Exemples fondamentaux}

\begin{exemple}[Catégories concrètes]\label{ex:categories-concretes}
\index{catégorie!concrète}
\begin{enumerate}[label=(\alph*)]
  \item $\Set$\index{Set@$\Set$} : objets = ensembles, morphismes = applications, composition usuelle.
  \item $\Grp$\index{Grp@$\Grp$} : objets = groupes, morphismes = homomorphismes de groupes.
  \item $\Ab$\index{Ab@$\Ab$} : objets = groupes abéliens, morphismes = homomorphismes de groupes.
  \item $\Ring$\index{Ring@$\Ring$} : objets = anneaux unitaires, morphismes = homomorphismes d'anneaux (préservant l'unité).
  \item $\Top$\index{Top@$\Top$} : objets = espaces topologiques, morphismes = applications continues.
  \item $\Vect_\K$\index{Vect@$\Vect_\K$} : objets = $\K$-espaces vectoriels, morphismes = applications $\K$-linéaires.
  \item $R$-$\Mod$\index{Mod@$R$-$\Mod$} : objets = $R$-modules à gauche, morphismes = homomorphismes de $R$-modules.
\end{enumerate}
\end{exemple}

\begin{exemple}[Ensembles ordonnés comme catégories]\label{ex:poset-categorie}
\index{catégorie!ensemble ordonné}
Soit $(P, \leq)$ un ensemble partiellement ordonné\index{ensemble ordonné}. On définit une catégorie $\mathcal{C}_P$ par :
\begin{itemize}
  \item $\Ob(\mathcal{C}_P) = P$ ;
  \item $\Hom_{\mathcal{C}_P}(x, y) = \begin{cases} \{(x,y)\} & \text{si } x \leq y, \\ \varnothing & \text{sinon.} \end{cases}$
\end{itemize}
La transitivité de $\leq$ assure l'existence de la composition, la réflexivité fournit les identités, et l'associativité est automatique (unicité des morphismes).
\end{exemple}

\begin{exemple}[Monoïdes comme catégories]\label{ex:monoide-categorie}
\index{catégorie!monoïde}\index{monoïde}
Tout monoïde $(M, \cdot, e)$ définit une catégorie $\mathcal{B}M$ à un seul objet $\star$ :
\[
  \Hom_{\mathcal{B}M}(\star, \star) = M, \quad g \circ f = g \cdot f, \quad \id_\star = e.
\]
Inversement, toute catégorie à un seul objet est (les données d'un) monoïde.
\end{exemple}

\begin{exemple}[Groupes comme catégories]\label{ex:groupe-categorie}
\index{catégorie!groupe}
Un groupe $G$ définit une catégorie $\mathcal{B}G$ à un seul objet où \emph{tous} les morphismes sont inversibles. Réciproquement, une catégorie à un seul objet dont tous les morphismes sont inversibles est (les données d')un groupe.
\end{exemple}

\subsection{Petitesse et taille}

\begin{definition}[Petite catégorie, catégorie localement petite]\label{def:petite-categorie}
\index{catégorie!petite}\index{catégorie!localement petite}
\begin{enumerate}[label=(\roman*)]
  \item Une catégorie $\mathcal{C}$ est \textbf{petite} si $\Ob(\mathcal{C})$ est un ensemble (et non une classe propre).
  \item Une catégorie $\mathcal{C}$ est \textbf{localement petite} si pour tous objets $A, B$, l'ensemble $\Hom_{\mathcal{C}}(A,B)$ est un ensemble (au sens de la théorie des ensembles fixée).
\end{enumerate}
\end{definition}

\begin{remarque}
Les catégories $\Set$, $\Grp$, $\Top$, etc.\ ne sont pas petites (la classe de tous les ensembles n'est pas un ensemble), mais elles sont localement petites. La catégorie des ensembles finis $\mathbf{FinSet}$ est un exemple de catégorie à la fois petite et localement petite.
\end{remarque}

\subsection{Monomorphismes, épimorphismes, isomorphismes}

\begin{definition}[Monomorphisme]\label{def:mono}
\index{monomorphisme}\index{morphisme!mono}
Un morphisme $f : A \to B$ dans $\mathcal{C}$ est un \textbf{monomorphisme} (ou \textbf{mono}) si pour tous morphismes $g, h : C \to A$,
\[
  f \circ g = f \circ h \implies g = h.
\]
On dit que $f$ est \emph{annulable à gauche}.
\end{definition}

\begin{definition}[Épimorphisme]\label{def:epi}
\index{épimorphisme}\index{morphisme!épi}
Un morphisme $f : A \to B$ dans $\mathcal{C}$ est un \textbf{épimorphisme} (ou \textbf{épi}) si pour tous morphismes $g, h : B \to C$,
\[
  g \circ f = h \circ f \implies g = h.
\]
On dit que $f$ est \emph{annulable à droite}.
\end{definition}

\begin{definition}[Isomorphisme]\label{def:iso}
\index{isomorphisme}\index{morphisme!iso}
Un morphisme $f : A \to B$ est un \textbf{isomorphisme} s'il existe $g : B \to A$ tel que
\[
  g \circ f = \id_A \quad \text{et} \quad f \circ g = \id_B.
\]
Un tel $g$ est unique et noté $f^{-1}$. On écrit $A \cong B$\index{$\cong$} si un tel isomorphisme existe.
\end{definition}

\begin{proposition}\label{prop:iso-mono-epi}
Tout isomorphisme est à la fois un monomorphisme et un épimorphisme.
\end{proposition}

\begin{proof}
Soit $f : A \to B$ un isomorphisme d'inverse $f^{-1}$. Si $f \circ g = f \circ h$, alors $g = f^{-1} \circ f \circ g = f^{-1} \circ f \circ h = h$, donc $f$ est mono. De même pour épi par composition à droite avec $f^{-1}$.
\end{proof}

\begin{remarque}
La réciproque est fausse en général. Dans $\Ring$, l'inclusion $\Z \hookrightarrow \Q$ est à la fois mono et épi, mais n'est pas un isomorphisme.
\end{remarque}

\begin{proposition}\label{prop:comp-mono-epi}
\begin{enumerate}[label=(\roman*)]
  \item La composée de deux monomorphismes est un monomorphisme.
  \item La composée de deux épimorphismes est un épimorphisme.
  \item Si $g \circ f$ est mono, alors $f$ est mono.
  \item Si $g \circ f$ est épi, alors $g$ est épi.
\end{enumerate}
\end{proposition}

\begin{proof}
(i) Soient $f : A \to B$ et $g : B \to C$ monos. Si $(g \circ f) \circ h_1 = (g \circ f) \circ h_2$, alors $g \circ (f \circ h_1) = g \circ (f \circ h_2)$, d'où $f \circ h_1 = f \circ h_2$ ($g$ mono), puis $h_1 = h_2$ ($f$ mono).

(iii) Si $g \circ f$ est mono et $f \circ h_1 = f \circ h_2$, alors $(g \circ f) \circ h_1 = (g \circ f) \circ h_2$, d'où $h_1 = h_2$. Les points (ii) et (iv) sont duaux.
\end{proof}

\begin{lemme}\label{lem:unicite-inverse}
\index{inverse!unicité}
L'inverse d'un isomorphisme est unique.
\end{lemme}

\begin{proof}
Si $g$ et $g'$ sont deux inverses de $f$, alors $g = g \circ \id_B = g \circ (f \circ g') = (g \circ f) \circ g' = \id_A \circ g' = g'$.
\end{proof}

\subsection{Objets initiaux, terminaux, nuls}

\begin{definition}[Objet initial, terminal, nul]\label{def:init-term-nul}
\index{objet!initial}\index{objet!terminal}\index{objet!nul}
Soit $\mathcal{C}$ une catégorie.
\begin{enumerate}[label=(\roman*)]
  \item Un objet $I$ est \textbf{initial} si pour tout objet $A$, il existe un unique morphisme $I \to A$.
  \item Un objet $T$ est \textbf{terminal} si pour tout objet $A$, il existe un unique morphisme $A \to T$.
  \item Un objet est \textbf{nul} (ou \textbf{zéro}) s'il est à la fois initial et terminal.
\end{enumerate}
\end{definition}

\begin{proposition}\label{prop:init-term-unique}
Si $I_1$ et $I_2$ sont deux objets initiaux de $\mathcal{C}$, alors $I_1 \cong I_2$, et l'isomorphisme est unique. De même pour les objets terminaux.
\end{proposition}

\begin{proof}
Par la propriété universelle, il existe des morphismes uniques $f : I_1 \to I_2$ et $g : I_2 \to I_1$. Alors $g \circ f : I_1 \to I_1$ et $\id_{I_1} : I_1 \to I_1$ sont deux morphismes de $I_1$ vers lui-même ; par unicité (car $I_1$ est initial), $g \circ f = \id_{I_1}$. De même $f \circ g = \id_{I_2}$.
\end{proof}

\begin{exemple}[Objets initiaux et terminaux]\label{ex:init-term}
\index{objet!initial!exemples}\index{objet!terminal!exemples}
\begin{enumerate}[label=(\alph*)]
  \item Dans $\Set$ : $\varnothing$ est initial, tout singleton est terminal.
  \item Dans $\Grp$ et $\Ab$ : le groupe trivial $\{e\}$ est objet nul.
  \item Dans $\Ring$ : $\Z$ est initial, l'anneau nul $\{0\}$ est terminal.
  \item Dans $\Top$ : $\varnothing$ est initial, tout singleton (avec la topologie triviale) est terminal.
  \item Dans $\Vect_\K$ : l'espace nul $\{0\}$ est objet nul.
\end{enumerate}
\end{exemple}

% ============================================================
\section{Foncteurs}

\subsection{Définition et premiers exemples}

\begin{definition}[Foncteur (covariant)]\label{def:foncteur}
\index{foncteur!covariant}\index{foncteur!définition}
Soient $\mathcal{C}$ et $\mathcal{D}$ deux catégories. Un \textbf{foncteur} (covariant) $F : \mathcal{C} \to \mathcal{D}$ est la donnée de :
\begin{enumerate}[label=(\roman*)]
  \item une application $F : \Ob(\mathcal{C}) \to \Ob(\mathcal{D})$ ;
  \item pour tout couple $(A, B)$ d'objets de $\mathcal{C}$, une application
  \[
    F_{A,B} : \Hom_{\mathcal{C}}(A, B) \longrightarrow \Hom_{\mathcal{D}}(F(A), F(B))
  \]
  (on note simplement $F(f)$ pour $F_{A,B}(f)$) ;
\end{enumerate}
satisfaisant :
\begin{enumerate}[label=(F\arabic*)]
  \item $F(\id_A) = \id_{F(A)}$ pour tout objet $A$ de $\mathcal{C}$ ;
  \item $F(g \circ f) = F(g) \circ F(f)$ pour tous morphismes composables $f, g$ dans $\mathcal{C}$.
\end{enumerate}
\end{definition}

\begin{definition}[Foncteur contravariant]\label{def:foncteur-contra}
\index{foncteur!contravariant}
Un \textbf{foncteur contravariant} $F : \mathcal{C} \to \mathcal{D}$ est un foncteur $F : \mathcal{C}^{\op} \to \mathcal{D}$, c'est-à-dire qu'il \emph{renverse les flèches} :
\[
  F(g \circ f) = F(f) \circ F(g).
\]
\end{definition}

\begin{exemple}[Foncteurs classiques]\label{ex:foncteurs}
\index{foncteur!exemples}
\begin{enumerate}[label=(\alph*)]
  \item \textbf{Foncteur d'oubli}\index{foncteur!d'oubli} $U : \Grp \to \Set$ : à un groupe $G$ associe l'ensemble sous-jacent, à un homomorphisme l'application sous-jacente.
  \item \textbf{Foncteur libre}\index{foncteur!libre} $F : \Set \to \Grp$ : à un ensemble $X$ associe le groupe libre $F(X)$.
  \item \textbf{Foncteur abélianisation}\index{abélianisation} $\mathrm{ab} : \Grp \to \Ab$ : $G \mapsto G/[G,G]$.
  \item \textbf{Foncteur espace dual}\index{foncteur!dual} $(-)^* : \Vect_\K \to \Vect_\K^{\op}$ : $V \mapsto V^* = \Hom_\K(V, \K)$, contravariant.
  \item \textbf{Foncteur $\pi_1$}\index{foncteur!$\pi_1$} : $\Top_* \to \Grp$, le groupe fondamental.
  \item \textbf{Foncteur identité}\index{foncteur!identité} $\Id_\mathcal{C} : \mathcal{C} \to \mathcal{C}$ : $A \mapsto A$, $f \mapsto f$.
  \item \textbf{Foncteur constant}\index{foncteur!constant} $\Delta_D : \mathcal{C} \to \mathcal{D}$ : tout objet sur $D$, tout morphisme sur $\id_D$.
\end{enumerate}
\end{exemple}

\begin{proposition}\label{prop:foncteur-iso}
\index{foncteur!préserve les isomorphismes}
Tout foncteur préserve les isomorphismes : si $f$ est un isomorphisme dans $\mathcal{C}$, alors $F(f)$ est un isomorphisme dans $\mathcal{D}$, d'inverse $F(f^{-1})$.
\end{proposition}

\begin{proof}
$F(f) \circ F(f^{-1}) = F(f \circ f^{-1}) = F(\id_B) = \id_{F(B)}$ et de même pour l'autre composition.
\end{proof}

\subsection{Fidélité, plénitude, essentielle surjectivité}

\begin{definition}[Foncteur fidèle, plein, pleinement fidèle]\label{def:fidele-plein}
\index{foncteur!fidèle}\index{foncteur!plein}\index{foncteur!pleinement fidèle}
Un foncteur $F : \mathcal{C} \to \mathcal{D}$ est :
\begin{enumerate}[label=(\roman*)]
  \item \textbf{fidèle} si pour tous $A, B \in \Ob(\mathcal{C})$, l'application $F_{A,B} : \Hom_\mathcal{C}(A,B) \to \Hom_\mathcal{D}(F(A), F(B))$ est injective ;
  \item \textbf{plein} si chaque $F_{A,B}$ est surjective ;
  \item \textbf{pleinement fidèle} s'il est à la fois plein et fidèle (chaque $F_{A,B}$ est bijective).
\end{enumerate}
\end{definition}

\begin{definition}[Essentiellement surjectif]\label{def:ess-surj}
\index{foncteur!essentiellement surjectif}
Un foncteur $F : \mathcal{C} \to \mathcal{D}$ est \textbf{essentiellement surjectif} si pour tout objet $D \in \Ob(\mathcal{D})$, il existe $C \in \Ob(\mathcal{C})$ tel que $F(C) \cong D$.
\end{definition}

\begin{exemple}\label{ex:foncteur-fidele}
\begin{enumerate}[label=(\alph*)]
  \item Le foncteur d'oubli $U : \Grp \to \Set$ est fidèle mais pas plein.
  \item L'inclusion $\Ab \hookrightarrow \Grp$ est pleinement fidèle.
  \item Le foncteur $\mathrm{ab} : \Grp \to \Ab$ est plein et essentiellement surjectif (mais pas fidèle).
\end{enumerate}
\end{exemple}

\begin{proposition}\label{prop:plein-fidele-mono}
Un foncteur fidèle reflète les monomorphismes : si $F$ est fidèle et $F(f)$ est mono, alors $f$ est mono.
\end{proposition}

\begin{proof}
Si $f \circ g = f \circ h$, alors $F(f) \circ F(g) = F(f \circ g) = F(f \circ h) = F(f) \circ F(h)$, donc $F(g) = F(h)$ car $F(f)$ est mono. Par fidélité, $g = h$.
\end{proof}

\subsection{Composition de foncteurs}

\begin{proposition}\label{prop:comp-foncteurs}
\index{foncteur!composition}
Si $F : \mathcal{C} \to \mathcal{D}$ et $G : \mathcal{D} \to \mathcal{E}$ sont des foncteurs, alors $G \circ F : \mathcal{C} \to \mathcal{E}$ défini par $(G \circ F)(A) = G(F(A))$ et $(G \circ F)(f) = G(F(f))$ est un foncteur.
\end{proposition}

\begin{proof}
$(G \circ F)(\id_A) = G(F(\id_A)) = G(\id_{F(A)}) = \id_{G(F(A))}$ et $(G \circ F)(g \circ f) = G(F(g \circ f)) = G(F(g) \circ F(f)) = G(F(g)) \circ G(F(f))$.
\end{proof}

% ============================================================
\section{Transformations naturelles}

\subsection{Définition}

\begin{definition}[Transformation naturelle]\label{def:transfo-nat}
\index{transformation naturelle!définition}
Soient $F, G : \mathcal{C} \to \mathcal{D}$ deux foncteurs. Une \textbf{transformation naturelle} $\alpha : F \Rightarrow G$ est la donnée, pour chaque objet $A$ de $\mathcal{C}$, d'un morphisme
\[
  \alpha_A : F(A) \longrightarrow G(A)
\]
dans $\mathcal{D}$, appelé \textbf{composante}\index{composante} de $\alpha$ en $A$, tel que pour tout morphisme $f : A \to B$ dans $\mathcal{C}$, le diagramme suivant commute :
\[
\begin{tikzcd}
  F(A) \arrow[r, "\alpha_A"] \arrow[d, "F(f)"'] & G(A) \arrow[d, "G(f)"] \\
  F(B) \arrow[r, "\alpha_B"'] & G(B)
\end{tikzcd}
\]
Autrement dit : $\alpha_B \circ F(f) = G(f) \circ \alpha_A$ pour tout $f : A \to B$.
\end{definition}

\begin{definition}[Isomorphisme naturel]\label{def:iso-nat}
\index{isomorphisme naturel}\index{transformation naturelle!isomorphisme}
Une transformation naturelle $\alpha : F \Rightarrow G$ est un \textbf{isomorphisme naturel} si chaque composante $\alpha_A$ est un isomorphisme dans $\mathcal{D}$. On écrit alors $F \cong G$.
\end{definition}

\begin{exemple}[Double dual]\label{ex:double-dual}
\index{double dual}
Soit $\K$ un corps. Le morphisme canonique $\eta_V : V \to V^{**}$, $v \mapsto (\varphi \mapsto \varphi(v))$, définit une transformation naturelle $\eta : \Id_{\Vect_\K} \Rightarrow (-)^{**}$.

Pour $f : V \to W$ linéaire, la naturalité revient à $f^{**} \circ \eta_V = \eta_W \circ f$, ce qui se vérifie : pour $v \in V$ et $\psi \in W^*$,
\[
  (f^{**}(\eta_V(v)))(\psi) = \eta_V(v)(\psi \circ f) = (\psi \circ f)(v) = \psi(f(v)) = \eta_W(f(v))(\psi).
\]
En dimension finie, $\eta$ est un isomorphisme naturel.
\end{exemple}

\begin{exemple}[Déterminant]\label{ex:determinant-nat}
\index{déterminant}
Le déterminant $\det : \mathrm{GL}_n \Rightarrow (-)^\times$ est une transformation naturelle entre foncteurs $\Ring \to \Grp$, où $\mathrm{GL}_n(R)$ est le groupe des matrices inversibles et $(-)^\times$ associe le groupe des unités. La naturalité exprime que pour tout homomorphisme d'anneaux $\varphi : R \to S$, on a $\det \circ \mathrm{GL}_n(\varphi) = \varphi^\times \circ \det$.
\end{exemple}

\begin{exemple}[Abélianisation]\label{ex:abel-nat}
\index{abélianisation!transformation naturelle}
La projection canonique $\pi_G : G \twoheadrightarrow G/[G,G]$ définit une transformation naturelle $\pi : U \Rightarrow U \circ \mathrm{ab}$ (ici $U : \Grp \to \Set$ est l'oubli), ou plus précisément $\pi : \Id_\Grp \Rightarrow \iota \circ \mathrm{ab}$ où $\iota : \Ab \hookrightarrow \Grp$.
\end{exemple}

\subsection{Composition verticale et horizontale}

\begin{definition}[Composition verticale]\label{def:comp-verticale}
\index{composition!verticale}
Soient $\alpha : F \Rightarrow G$ et $\beta : G \Rightarrow H$ des transformations naturelles entre foncteurs $\mathcal{C} \to \mathcal{D}$. La \textbf{composition verticale} $\beta \circ \alpha : F \Rightarrow H$ est définie par
\[
  (\beta \circ \alpha)_A = \beta_A \circ \alpha_A.
\]
\[
\begin{tikzcd}
  F(A) \arrow[r, "\alpha_A"] \arrow[rr, bend left=40, "(\beta\circ\alpha)_A"] & G(A) \arrow[r, "\beta_A"] & H(A)
\end{tikzcd}
\]
\end{definition}

\begin{proposition}\label{prop:comp-vert-naturelle}
La composition verticale $\beta \circ \alpha$ est bien une transformation naturelle.
\end{proposition}

\begin{proof}
Pour $f : A \to B$ dans $\mathcal{C}$ :
$(\beta \circ \alpha)_B \circ F(f) = \beta_B \circ \alpha_B \circ F(f) = \beta_B \circ G(f) \circ \alpha_A = H(f) \circ \beta_A \circ \alpha_A = H(f) \circ (\beta \circ \alpha)_A$.
\end{proof}

\begin{definition}[Composition horizontale]\label{def:comp-horizontale}
\index{composition!horizontale}
Soient $F, G : \mathcal{C} \to \mathcal{D}$ et $H, K : \mathcal{D} \to \mathcal{E}$ des foncteurs, avec $\alpha : F \Rightarrow G$ et $\beta : H \Rightarrow K$. La \textbf{composition horizontale} $\beta * \alpha : H \circ F \Rightarrow K \circ G$ est définie par
\[
  (\beta * \alpha)_A = \beta_{G(A)} \circ H(\alpha_A) = K(\alpha_A) \circ \beta_{F(A)}.
\]
\[
\begin{tikzcd}
  \mathcal{C}
    \arrow[r, bend left=30, "F", ""{name=U,below}]
    \arrow[r, bend right=30, "G"', ""{name=D}]
  & \mathcal{D}
    \arrow[r, bend left=30, "H", ""{name=V,below}]
    \arrow[r, bend right=30, "K"', ""{name=W}]
  & \mathcal{E}
  \arrow[Rightarrow, from=U, to=D, "\alpha"]
  \arrow[Rightarrow, from=V, to=W, "\beta"]
\end{tikzcd}
\]
\end{definition}

\begin{proposition}[Loi d'échange]\label{prop:loi-echange}
\index{loi d'échange}
Soient $\alpha : F \Rightarrow G$, $\alpha' : G \Rightarrow H$ (foncteurs $\mathcal{C} \to \mathcal{D}$) et $\beta : J \Rightarrow K$, $\beta' : K \Rightarrow L$ (foncteurs $\mathcal{D} \to \mathcal{E}$). Alors
\[
  (\beta' \circ \beta) * (\alpha' \circ \alpha) = (\beta' * \alpha') \circ (\beta * \alpha).
\]
\end{proposition}

\begin{proof}
Vérification composante par composante en utilisant la naturalité de $\beta$ et $\beta'$.
\end{proof}

\subsection{Catégorie de foncteurs}

\begin{definition}[Catégorie de foncteurs]\label{def:cat-foncteurs}
\index{catégorie!de foncteurs}\index{$[\mathcal{C},\mathcal{D}]$}
Soient $\mathcal{C}$ et $\mathcal{D}$ deux catégories (avec $\mathcal{C}$ petite). La \textbf{catégorie de foncteurs} $[\mathcal{C}, \mathcal{D}]$ (aussi notée $\Fun(\mathcal{C}, \mathcal{D})$ ou $\mathcal{D}^\mathcal{C}$) est définie par :
\begin{itemize}
  \item les objets sont les foncteurs $F : \mathcal{C} \to \mathcal{D}$ ;
  \item les morphismes sont les transformations naturelles $\alpha : F \Rightarrow G$ ;
  \item la composition est la composition verticale ;
  \item l'identité de $F$ est la transformation naturelle $\id_F$ définie par $(\id_F)_A = \id_{F(A)}$.
\end{itemize}
\end{definition}

\begin{remarque}
L'hypothèse « $\mathcal{C}$ petite » est nécessaire pour garantir que les transformations naturelles entre deux foncteurs forment un ensemble (et non une classe propre).
\end{remarque}

\begin{proposition}\label{prop:iso-cat-foncteurs}
Un morphisme $\alpha : F \to G$ dans $[\mathcal{C}, \mathcal{D}]$ est un isomorphisme si et seulement si $\alpha$ est un isomorphisme naturel (i.e.\ chaque $\alpha_A$ est un isomorphisme dans $\mathcal{D}$).
\end{proposition}

\begin{proof}
Si $\alpha$ est un isomorphisme dans $[\mathcal{C}, \mathcal{D}]$, il existe $\beta : G \Rightarrow F$ avec $\beta \circ \alpha = \id_F$ et $\alpha \circ \beta = \id_G$. En composantes : $\beta_A \circ \alpha_A = \id_{F(A)}$ et $\alpha_A \circ \beta_A = \id_{G(A)}$, donc chaque $\alpha_A$ est un isomorphisme.

Réciproquement, si chaque $\alpha_A$ est un isomorphisme, on pose $\beta_A = \alpha_A^{-1}$. La naturalité de $\beta$ découle de celle de $\alpha$ : de $\alpha_B \circ F(f) = G(f) \circ \alpha_A$, on tire $F(f) = \alpha_B^{-1} \circ G(f) \circ \alpha_A$, d'où $\beta_B \circ G(f) = F(f) \circ \beta_A$.
\end{proof}

\begin{exemple}[Préfaisceaux]\label{ex:prefaisceaux}
\index{préfaisceau}
Si $\mathcal{C}$ est une petite catégorie, un \textbf{préfaisceau} (d'ensembles) sur $\mathcal{C}$ est un foncteur $F : \mathcal{C}^{\op} \to \Set$. La catégorie des préfaisceaux est $[\mathcal{C}^{\op}, \Set]$, souvent notée $\widehat{\mathcal{C}}$.
\end{exemple}

% ============================================================
\section{Exercices}

\begin{exercice}[Catégorie des flèches]\label{exo:cat-fleches}
\index{catégorie!des flèches}
Soit $\mathcal{C}$ une catégorie. Construire la \textbf{catégorie des flèches} $\mathcal{C}^{\to}$ dont les objets sont les morphismes de $\mathcal{C}$ et dont un morphisme de $f : A \to B$ vers $g : C \to D$ est un couple $(h, k)$ tel que le diagramme suivant commute :
\[
\begin{tikzcd}
  A \arrow[r, "h"] \arrow[d, "f"'] & C \arrow[d, "g"] \\
  B \arrow[r, "k"'] & D
\end{tikzcd}
\]
Montrer que $\mathcal{C}^{\to} \cong [\mathbf{2}, \mathcal{C}]$ où $\mathbf{2}$ est la catégorie $\{0 \to 1\}$.
\end{exercice}

\begin{exercice}[Sections et rétractions]\label{exo:section-retraction}
\index{section}\index{rétraction}
Un morphisme $f : A \to B$ est une \textbf{section} (ou mono scindé) s'il existe $g : B \to A$ tel que $g \circ f = \id_A$. Il est une \textbf{rétraction} (ou épi scindé) s'il existe $h : B \to A$ tel que $f \circ h = \id_B$.
\begin{enumerate}[label=(\alph*)]
  \item Montrer que toute section est un monomorphisme et toute rétraction est un épimorphisme.
  \item Donner un monomorphisme dans $\Set$ qui n'est pas une section.
  \item Montrer que dans $\Set$, tout épimorphisme est une rétraction (utiliser l'axiome du choix).
\end{enumerate}
\end{exercice}

\begin{exercice}[Foncteur puissance]\label{exo:foncteur-puissance}
\index{foncteur!puissance}
Montrer que l'application $\mathcal{P} : \Set \to \Set$ qui à un ensemble $X$ associe $\mathcal{P}(X)$ (ensemble des parties) peut être munie d'une structure de foncteur de deux façons :
\begin{enumerate}[label=(\alph*)]
  \item Covariant : image directe $f_*(S) = f(S)$.
  \item Contravariant : image réciproque $f^*(S) = f^{-1}(S)$.
\end{enumerate}
Vérifier les axiomes de foncteur dans chaque cas.
\end{exercice}

\begin{exercice}[Naturalité du centre]\label{exo:centre}
\index{centre d'un groupe}
Le centre $Z(G) = \{g \in G \mid \forall h \in G,\, gh = hg\}$ définit-il un foncteur $Z : \Grp \to \Ab$ ? Si oui, l'inclusion $Z(G) \hookrightarrow G$ est-elle une transformation naturelle ? Justifier.
\end{exercice}

\begin{exercice}[Équivalence mono/injectif dans Set]\label{exo:mono-set}
Montrer que dans $\Set$, un morphisme est un monomorphisme si et seulement s'il est injectif, et un épimorphisme si et seulement s'il est surjectif.
\end{exercice}

\begin{exercice}[Catégories de tranches]\label{exo:tranche}
\index{catégorie!de tranches}
Soit $\mathcal{C}$ une catégorie et $A$ un objet. Construire la \textbf{catégorie de tranches} $\mathcal{C}/A$ (ou \emph{comma category over $A$}) dont les objets sont les morphismes $f : X \to A$ et les morphismes de $(f : X \to A)$ vers $(g : Y \to A)$ sont les morphismes $h : X \to Y$ tels que $g \circ h = f$. Vérifier les axiomes de catégorie. Décrire $\Set/\{0,1\}$.
\end{exercice}

\begin{exercice}[Endomorphismes de l'identité]\label{exo:endo-id}
\index{transformation naturelle!endomorphismes de l'identité}
Soit $\mathcal{C} = \Ab$. Montrer que $\End({\Id_\Ab}) \cong \Z$ (l'anneau des endomorphismes de la transformation naturelle identité est isomorphe à $\Z$).

\emph{Indication} : une transformation naturelle $\alpha : \Id_\Ab \Rightarrow \Id_\Ab$ est déterminée par $\alpha_\Z(1)$.
\end{exercice}

\begin{exercice}[Foncteur nerf]\label{exo:nerf}
\index{foncteur!nerf}
Soit $\mathcal{C}$ une petite catégorie. Le \textbf{nerf}\index{nerf} $N(\mathcal{C})$ est l'ensemble simplicial défini par $N(\mathcal{C})_n = \{$chaînes de $n$ morphismes composables $A_0 \xrightarrow{f_1} A_1 \xrightarrow{f_2} \cdots \xrightarrow{f_n} A_n\}$.
\begin{enumerate}[label=(\alph*)]
  \item Décrire $N(\mathcal{C})_0$, $N(\mathcal{C})_1$ et $N(\mathcal{C})_2$.
  \item Montrer que $\mathcal{C} \mapsto N(\mathcal{C})$ définit un foncteur $N : \Cat \to \mathbf{sSet}$.
\end{enumerate}
\end{exercice}


% ============================================================
% CHAPITRE 2
% ============================================================
\chapter{Dualité et Principe de Dualité}
\minitoc

\section{La catégorie opposée}

\begin{definition}[Catégorie opposée]\label{def:cat-opposee}
\index{catégorie!opposée}\index{$\mathcal{C}^{\op}$}
Soit $\mathcal{C}$ une catégorie. La \textbf{catégorie opposée} (ou \textbf{duale}) $\mathcal{C}^{\op}$ est définie par :
\begin{enumerate}[label=(\roman*)]
  \item $\Ob(\mathcal{C}^{\op}) = \Ob(\mathcal{C})$ ;
  \item $\Hom_{\mathcal{C}^{\op}}(A, B) = \Hom_{\mathcal{C}}(B, A)$ ;
  \item la composition dans $\mathcal{C}^{\op}$ : si $f^{\op} \in \Hom_{\mathcal{C}^{\op}}(A,B)$ et $g^{\op} \in \Hom_{\mathcal{C}^{\op}}(B,C)$ (correspondant à $f : B \to A$ et $g : C \to B$ dans $\mathcal{C}$), alors $g^{\op} \circ_{\op} f^{\op} = (f \circ g)^{\op}$ ;
  \item $\id_A^{\op} = (\id_A)^{\op}$.
\end{enumerate}
\end{definition}

\begin{proposition}\label{prop:cop-cop}
Pour toute catégorie $\mathcal{C}$, on a $(\mathcal{C}^{\op})^{\op} = \mathcal{C}$.
\end{proposition}

\begin{proof}
Les objets sont identiques. Les morphismes : $\Hom_{(\mathcal{C}^{\op})^{\op}}(A,B) = \Hom_{\mathcal{C}^{\op}}(B,A) = \Hom_\mathcal{C}(A,B)$. La composition est restaurée par double renversement.
\end{proof}

\begin{exemple}\label{ex:cop-exemples}
\begin{enumerate}[label=(\alph*)]
  \item Si $(P, \leq)$ est un ensemble ordonné, $\mathcal{C}_P^{\op} = \mathcal{C}_{(P, \geq)}$.
  \item Si $G$ est un groupe vu comme catégorie à un objet, $(\mathcal{B}G)^{\op} = \mathcal{B}G^{\op}$ où $G^{\op}$ est le groupe opposé (même ensemble, multiplication $a \cdot^{\op} b = b \cdot a$). Pour les groupes abéliens, $G^{\op} \cong G$.
  \item $\Vect_\K^{\op}$ n'est pas « concrètement » la catégorie des espaces vectoriels, mais elle lui est équivalente (via le dual).
\end{enumerate}
\end{exemple}

\section{Le principe de dualité}

\begin{theoreme}[Principe de dualité]\label{thm:dualite}
\index{dualité!principe de}
Soit $\Phi$ un énoncé formulé dans le langage de la théorie des catégories (en termes d'objets, morphismes, composition et identités). Soit $\Phi^{\op}$ l'énoncé obtenu en renversant toutes les flèches (en échangeant source et but, en inversant l'ordre de composition). Alors :
\begin{center}
  $\Phi$ est vrai dans toute catégorie $\iff$ $\Phi^{\op}$ est vrai dans toute catégorie.
\end{center}
Plus précisément, $\Phi$ est vrai dans $\mathcal{C}$ si et seulement si $\Phi^{\op}$ est vrai dans $\mathcal{C}^{\op}$.
\end{theoreme}

\begin{proof}
Si $\Phi$ est vrai dans toute catégorie, alors en particulier $\Phi$ est vrai dans $\mathcal{C}^{\op}$ pour toute catégorie $\mathcal{C}$. Or l'énoncé « $\Phi$ dans $\mathcal{C}^{\op}$ » est exactement « $\Phi^{\op}$ dans $\mathcal{C}$ » (car passer à l'opposée renverse les flèches). Donc $\Phi^{\op}$ est vrai dans toute catégorie $\mathcal{C}$.
\end{proof}

\begin{remarque}
Le principe de dualité est un \emph{méta-théorème} : il porte sur les énoncés de la théorie et non sur des objets mathématiques particuliers. Son utilité pratique est considérable : tout théorème catégorique donne \emph{gratuitement} un second théorème par dualisation.
\end{remarque}

\begin{corollaire}\label{cor:dual-concepts}
Voici un dictionnaire des concepts duaux :
\begin{center}
\begin{tabular}{ll}
\toprule
\textbf{Concept} & \textbf{Concept dual} \\
\midrule
monomorphisme & épimorphisme \\
objet initial & objet terminal \\
section (mono scindé) & rétraction (épi scindé) \\
noyau & conoyau \\
produit & coproduit \\
limite & colimite \\
égaliseur & coégaliseur \\
\bottomrule
\end{tabular}
\end{center}
\end{corollaire}

\section{Exemples de dualisation}

\begin{exemple}[Dualisation des propositions du chapitre 1]\label{ex:dual-chap1}
\index{dualité!exemples}
\begin{enumerate}[label=(\alph*)]
  \item « La composée de monomorphismes est un monomorphisme » dualise en « la composée d'épimorphismes est un épimorphisme ».
  \item « Si $g \circ f$ est mono alors $f$ est mono » dualise en « si $g \circ f$ est épi alors $g$ est épi ».
  \item « L'objet initial est unique à isomorphisme unique près » dualise en la même affirmation pour l'objet terminal.
\end{enumerate}
\end{exemple}

\begin{proposition}\label{prop:contra-cop}
\index{foncteur!contravariant!et catégorie opposée}
Un foncteur contravariant $F : \mathcal{C} \to \mathcal{D}$ est la même chose qu'un foncteur (covariant) $F : \mathcal{C}^{\op} \to \mathcal{D}$, ou de façon équivalente un foncteur $F : \mathcal{C} \to \mathcal{D}^{\op}$.
\end{proposition}

\begin{proof}
La donnée d'une application $F : \Hom_\mathcal{C}(A,B) \to \Hom_\mathcal{D}(F(B), F(A))$ vérifiant $F(g \circ f) = F(f) \circ F(g)$ est exactement la donnée d'une application $\Hom_{\mathcal{C}^{\op}}(B,A) \to \Hom_\mathcal{D}(F(B), F(A))$ vérifiant la fonctorialité covariante (en utilisant la composition dans $\mathcal{C}^{\op}$).
\end{proof}

\section{Dualités concrètes}

Les dualités abstraites de la théorie des catégories se manifestent dans des théorèmes classiques profonds.

\subsection{Dualité de Pontryagin}

\begin{theoreme}[Pontryagin]\label{thm:pontryagin}
\index{dualité!de Pontryagin}\index{Pontryagin}
Soit $\mathbf{LCA}$ la catégorie des groupes abéliens localement compacts et homomorphismes continus. Le foncteur \textbf{dual de Pontryagin}
\[
  \widehat{(-)} : \mathbf{LCA} \longrightarrow \mathbf{LCA}^{\op}, \quad G \longmapsto \widehat{G} = \Hom_{\mathrm{cont}}(G, \R/\Z)
\]
(muni de la topologie compacte-ouverte) définit une équivalence de catégories $\mathbf{LCA} \simeq \mathbf{LCA}^{\op}$.

En particulier, la transformation naturelle canonique $G \to \widehat{\widehat{G}}$ est un isomorphisme naturel.
\end{theoreme}

\begin{remarque}
La dualité de Pontryagin échange :
\begin{itemize}
  \item groupes compacts $\longleftrightarrow$ groupes discrets ;
  \item $\Z \longleftrightarrow \R/\Z$ (le cercle) ;
  \item $\Z/n\Z \longleftrightarrow \Z/n\Z$ (auto-dual).
\end{itemize}
\end{remarque}

\subsection{Dualité de Stone}

\begin{theoreme}[Stone]\label{thm:stone}
\index{dualité!de Stone}\index{Stone}
Il existe une équivalence de catégories
\[
  \mathbf{Bool}^{\op} \simeq \mathbf{Stone}
\]
où $\mathbf{Bool}$ est la catégorie des algèbres de Boole et $\mathbf{Stone}$ est la catégorie des espaces de Stone (espaces topologiques compacts totalement discontinus).

Le foncteur $\mathbf{Bool} \to \mathbf{Stone}^{\op}$ associe à une algèbre de Boole $B$ son espace de spectres $\mathrm{Spec}(B)$ (ensemble des ultrafiltres muni de la topologie de Zariski). Le foncteur inverse associe à un espace de Stone $X$ l'algèbre $\mathrm{Clop}(X)$ de ses parties clopen (ouvertes et fermées).
\end{theoreme}

\begin{remarque}
La dualité de Stone peut être vue comme un précurseur de la dualité de Gelfand (entre $C^*$-algèbres commutatives et espaces compacts) et plus généralement des dualités en géométrie algébrique (le foncteur $\mathrm{Spec}$).
\end{remarque}

\section{Le foncteur Hom}

\begin{definition}[Foncteur Hom covariant]\label{def:hom-covariant}
\index{foncteur!Hom!covariant}\index{Hom@$\Hom$}
Soit $\mathcal{C}$ une catégorie localement petite et $A \in \Ob(\mathcal{C})$. Le \textbf{foncteur Hom covariant} est
\[
  \Hom_\mathcal{C}(A, -) : \mathcal{C} \longrightarrow \Set
\]
défini par :
\begin{itemize}
  \item sur les objets : $B \mapsto \Hom_\mathcal{C}(A, B)$ ;
  \item sur les morphismes : à $f : B \to C$, on associe la post-composition
  \[
    f_* = \Hom_\mathcal{C}(A, f) : \Hom_\mathcal{C}(A, B) \to \Hom_\mathcal{C}(A, C), \quad g \mapsto f \circ g.
  \]
\end{itemize}
\end{definition}

\begin{proposition}\label{prop:hom-cov-foncteur}
$\Hom_\mathcal{C}(A, -)$ est bien un foncteur.
\end{proposition}

\begin{proof}
$(\id_B)_*(g) = \id_B \circ g = g$, donc $(\id_B)_* = \id_{\Hom(A,B)}$. Pour la composition : $(f' \circ f)_*(g) = (f' \circ f) \circ g = f' \circ (f \circ g) = f'_*(f_*(g))$.
\end{proof}

\begin{definition}[Foncteur Hom contravariant]\label{def:hom-contravariant}
\index{foncteur!Hom!contravariant}
Soit $B \in \Ob(\mathcal{C})$. Le \textbf{foncteur Hom contravariant} est
\[
  \Hom_\mathcal{C}(-, B) : \mathcal{C}^{\op} \longrightarrow \Set
\]
défini par :
\begin{itemize}
  \item sur les objets : $A \mapsto \Hom_\mathcal{C}(A, B)$ ;
  \item sur les morphismes : à $f : A' \to A$ dans $\mathcal{C}$, on associe la pré-composition
  \[
    f^* = \Hom_\mathcal{C}(f, B) : \Hom_\mathcal{C}(A, B) \to \Hom_\mathcal{C}(A', B), \quad g \mapsto g \circ f.
  \]
\end{itemize}
\end{definition}

\begin{definition}[Bifoncteur Hom]\label{def:hom-bifoncteur}
\index{foncteur!Hom!bifoncteur}\index{bifoncteur}
Le \textbf{bifoncteur Hom} est le foncteur
\[
  \Hom_\mathcal{C}(-, -) : \mathcal{C}^{\op} \times \mathcal{C} \longrightarrow \Set
\]
défini par $(A, B) \mapsto \Hom_\mathcal{C}(A, B)$ et $(f, g) \mapsto (h \mapsto g \circ h \circ f)$ pour $f : A' \to A$ et $g : B \to B'$.
\end{definition}

\begin{proposition}\label{prop:hom-bifoncteur}
$\Hom_\mathcal{C}(-, -)$ est bien un foncteur $\mathcal{C}^{\op} \times \mathcal{C} \to \Set$.
\end{proposition}

\begin{proof}
Il faut vérifier les deux axiomes. Pour l'identité : $\Hom(\id_A, \id_B)(h) = \id_B \circ h \circ \id_A = h$. Pour la composition : si $(f_1, g_1)$ et $(f_2, g_2)$ sont composables, alors
\begin{align*}
  \Hom(f_1 \circ f_2, g_2 \circ g_1)(h) &= (g_2 \circ g_1) \circ h \circ (f_1 \circ f_2) \\
  &= g_2 \circ (g_1 \circ h \circ f_1) \circ f_2 \\
  &= \Hom(f_2, g_2)(\Hom(f_1, g_1)(h)).
\end{align*}
(Attention à l'ordre dans $\mathcal{C}^{\op}$ : la composition de $(f_2^{\op}, g_2)$ après $(f_1^{\op}, g_1)$ en tant que morphisme de $\mathcal{C}^{\op} \times \mathcal{C}$ donne $((f_1 \circ f_2)^{\op}, g_2 \circ g_1)$.)
\end{proof}

\begin{proposition}\label{prop:hom-mono-epi}
\index{monomorphisme!caractérisation par Hom}\index{épimorphisme!caractérisation par Hom}
Soit $\mathcal{C}$ localement petite et $f : A \to B$ un morphisme.
\begin{enumerate}[label=(\roman*)]
  \item $f$ est mono $\iff$ $f_* : \Hom(C, A) \to \Hom(C, B)$ est injective pour tout objet $C$.
  \item $f$ est épi $\iff$ $f^* : \Hom(B, C) \to \Hom(A, C)$ est injective pour tout objet $C$.
\end{enumerate}
\end{proposition}

\begin{proof}
Ce sont exactement les reformulations des définitions de mono et épi en termes des applications induites.
\end{proof}

\begin{proposition}\label{prop:hom-transfo-nat}
\index{transformation naturelle!entre foncteurs Hom}
Un morphisme $f : A \to B$ dans $\mathcal{C}$ induit une transformation naturelle
\[
  f^* : \Hom_\mathcal{C}(B, -) \Longrightarrow \Hom_\mathcal{C}(A, -)
\]
définie par $f^*_C(g) = g \circ f$ pour $g : B \to C$.
\end{proposition}

\begin{proof}
Pour $h : C \to D$, on vérifie la naturalité : $h_* \circ f^*_C = f^*_D \circ h_*$. En effet, pour $g : B \to C$ :
$(h_* \circ f^*_C)(g) = h \circ (g \circ f) = (h \circ g) \circ f = f^*_D(h \circ g) = (f^*_D \circ h_*)(g)$.
\end{proof}

\section{Exercices}

\begin{exercice}[Catégorie opposée de la catégorie opposée]\label{exo:cop-cop}
Vérifier en détail que $(\mathcal{C}^{\op})^{\op} = \mathcal{C}$ (égalité et non simple isomorphisme).
\end{exercice}

\begin{exercice}[Dual d'un foncteur]\label{exo:dual-foncteur}
Soit $F : \mathcal{C} \to \mathcal{D}$ un foncteur. Montrer que $F$ induit un foncteur $F^{\op} : \mathcal{C}^{\op} \to \mathcal{D}^{\op}$ défini par $F^{\op}(A) = F(A)$ et $F^{\op}(f^{\op}) = (F(f))^{\op}$. Vérifier les axiomes.
\end{exercice}

\begin{exercice}[Pleinement fidèle et isomorphismes]\label{exo:pf-iso}
\index{foncteur!pleinement fidèle!et isomorphismes}
Soit $F : \mathcal{C} \to \mathcal{D}$ un foncteur pleinement fidèle. Montrer que $F$ reflète les isomorphismes : si $F(f)$ est un isomorphisme dans $\mathcal{D}$, alors $f$ est un isomorphisme dans $\mathcal{C}$.
\end{exercice}

\begin{exercice}[Caractérisation des mono dans $\Top$]\label{exo:mono-top}
\index{monomorphisme!dans $\Top$}
\begin{enumerate}[label=(\alph*)]
  \item Montrer que dans $\Top$, les monomorphismes sont exactement les applications continues injectives.
  \item Montrer que dans $\Top$, les épimorphismes sont exactement les applications continues surjectives.
  \item En déduire qu'il existe dans $\Top$ des morphismes qui sont à la fois mono et épi sans être des isomorphismes. Donner un exemple explicite.
\end{enumerate}
\end{exercice}

\begin{exercice}[Hom et produit]\label{exo:hom-produit}
Soit $\mathcal{C}$ une catégorie localement petite possédant des produits binaires. Montrer qu'il existe un isomorphisme naturel (en $C$)
\[
  \Hom_\mathcal{C}(C, A \times B) \cong \Hom_\mathcal{C}(C, A) \times \Hom_\mathcal{C}(C, B).
\]
\end{exercice}

\begin{exercice}[Dualité de Gelfand -- énoncé catégorique]\label{exo:gelfand}
\index{dualité!de Gelfand}
Soit $\mathbf{CHaus}$ la catégorie des espaces compacts de Hausdorff et $\mathbf{cC^*Alg}$ la catégorie des $C^*$-algèbres commutatives unitaires.
\begin{enumerate}[label=(\alph*)]
  \item Décrire le foncteur $C : \mathbf{CHaus} \to \mathbf{cC^*Alg}^{\op}$ qui à $X$ associe $C(X, \C)$.
  \item Énoncer le théorème de Gelfand--Naimark en termes d'équivalence de catégories.
  \item En quoi cette dualité généralise-t-elle la dualité de Stone ?
\end{enumerate}
\end{exercice}

\begin{exercice}[Foncteur Hom interne]\label{exo:hom-interne}
\index{Hom interne}
Dans $\Vect_\K$, montrer que $\Hom_\K(V, W)$ est lui-même un $\K$-espace vectoriel, de sorte que le bifoncteur $\Hom$ se factorise par $\Vect_\K$ plutôt que par $\Set$. On dit que $\Vect_\K$ est \textbf{enrichie sur elle-même}. Est-ce le cas pour $\Grp$ ?
\end{exercice}

\begin{exercice}[Dualité dans les ensembles ordonnés]\label{exo:dual-poset}
Soit $(P, \leq)$ un ensemble ordonné fini.
\begin{enumerate}[label=(\alph*)]
  \item Montrer que les éléments minimaux de $P$ correspondent aux objets initiaux de $\mathcal{C}_P$, et les éléments maximaux aux objets terminaux.
  \item En déduire que la catégorie $\mathcal{C}_P$ admet un objet initial si et seulement si $P$ a un plus petit élément.
  \item Appliquer le principe de dualité pour obtenir le résultat analogue pour les objets terminaux.
\end{enumerate}
\end{exercice}

\begin{exercice}[Plongement de Yoneda -- première rencontre]\label{exo:yoneda-intro}
\index{Yoneda!plongement}
Soit $\mathcal{C}$ une catégorie localement petite. On définit le \textbf{plongement de Yoneda}\index{plongement de Yoneda}
\[
  \mathsf{y} : \mathcal{C} \longrightarrow [\mathcal{C}^{\op}, \Set], \quad A \longmapsto \Hom_\mathcal{C}(-, A).
\]
\begin{enumerate}[label=(\alph*)]
  \item Montrer que $\mathsf{y}$ est un foncteur (à un morphisme $f : A \to B$, associer la transformation naturelle $f_* : \Hom(-, A) \Rightarrow \Hom(-, B)$).
  \item Montrer que $\mathsf{y}$ est pleinement fidèle. (C'est une partie du \textbf{lemme de Yoneda}, qui sera traité en détail au chapitre suivant.)
\end{enumerate}
\end{exercice}
